\chapter{Probability Models and Rules}

\section{The axioms of probability}

Chapter 1 assigned probabilities to outcomes of a discrete sample space. To handle all
sample spaces (including continuous ones) uniformly, we take the following properties as
axioms and derive everything else from them.

\begin{definition}[Probability model]
A \emph{probability model} consists of a sample space $S$, a collection of events
(subsets of $S$), and a function $P$ assigning a real number $P(A)$ to each event, such
that
\begin{enumerate}
  \item[\textbf{A1.}] $P(S) = 1$;
  \item[\textbf{A2.}] $P(A) \ge 0$ for every event $A$;
  \item[\textbf{A3.}] (countable additivity) if $A_1, A_2, \dots$ are pairwise mutually
        exclusive, then
        \[
          P\Big(\bigcup_{i=1}^\infty A_i\Big) = \sum_{i=1}^\infty P(A_i).
        \]
\end{enumerate}
\end{definition}

Taking $A_{k+1} = A_{k+2} = \dots = \emptyset$ in A3 (once we know $P(\emptyset)=0$, see
below) gives \emph{finite additivity}: for pairwise disjoint $A_1,\dots,A_k$,
$P(A_1 \cup \dots \cup A_k) = P(A_1) + \dots + P(A_k)$.

\begin{remark}
The discrete assignment of Chapter 1 ($P(A) = \sum_{a\in A} P(a)$) satisfies all three
axioms, so it is a special case of this definition.
\end{remark}

\section{Consequences of the axioms}

\begin{theorem}[Basic rules]
For any events $A$, $B$:
\begin{enumerate}
  \item $P(\emptyset) = 0$.
  \item (Complement rule) $P(\bar A) = 1 - P(A)$.
  \item If $A \subseteq B$ then $P(A) \le P(B)$.
  \item $0 \le P(A) \le 1$.
\end{enumerate}
\end{theorem}

\begin{proof}
(1) Apply A3 to $A_1 = A_2 = \dots = \emptyset$ (pairwise disjoint): $P(\emptyset) =
\sum_{i} P(\emptyset)$, which forces $P(\emptyset) = 0$ since $P(\emptyset)$ is a finite
non-negative number.

(2) $A$ and $\bar A$ are disjoint with union $S$, so $1 = P(S) = P(A) + P(\bar A)$.

(3) Write $B = A \cup (B\cap \bar A)$, a disjoint union. Then
$P(B) = P(A) + P(B \cap \bar A) \ge P(A)$ by A2.

(4) $P(A)\ge 0$ is A2, and $A \subseteq S$ gives $P(A) \le P(S) = 1$ by (3).
\end{proof}

The complement rule is one of the most useful tools in the course: when ``at least one''
appears, compute the probability of ``none'' instead.

\begin{example}[At least one six]
Roll a fair die four times. The outcomes (ordered $4$-tuples) are $6^4$ equally likely
possibilities, of which $5^4$ contain no six. Hence
\[
  P(\text{at least one six}) = 1 - \frac{5^4}{6^4} = 1 - \frac{625}{1296} = \frac{671}{1296}
  \approx 0.5177.
\]
By contrast, for $24$ rolls of a pair of dice, $P(\text{at least one double six}) =
1 - (35/36)^{24} \approx 0.4914$. (This comparison is the famous Chevalier de M\'er\'e
problem: he noticed by experience that the second bet was slightly unfavourable.)
\end{example}

\section{Inclusion--exclusion}

\begin{theorem}[Inclusion--exclusion for two events]
For any events $A$ and $B$,
\[
  P(A \cup B) = P(A) + P(B) - P(A\cap B).
\]
\end{theorem}

\begin{proof}
$A \cup B$ is the disjoint union of $A$ and $\bar A \cap B$, so
$P(A\cup B) = P(A) + P(\bar A\cap B)$. Also $B$ is the disjoint union of $A\cap B$ and
$\bar A \cap B$, so $P(\bar A\cap B) = P(B) - P(A\cap B)$. Substitute.
\end{proof}

Intuitively, adding $P(A)$ and $P(B)$ counts the overlap $A\cap B$ twice, so we subtract
it once.

\begin{theorem}[Inclusion--exclusion for three events]
\[
  P(A\cup B\cup C) = P(A)+P(B)+P(C) - P(AB) - P(AC) - P(BC) + P(ABC).
\]
More generally, for $A_1,\dots,A_n$,
\[
  P\Big(\bigcup_{i=1}^n A_i\Big) = \sum_i P(A_i) - \sum_{i<j} P(A_iA_j)
  + \sum_{i<j<k} P(A_iA_jA_k) - \dots + (-1)^{n+1} P(A_1A_2\cdots A_n).
\]
\end{theorem}

\begin{proof}[Proof sketch]
Apply the two-event rule to $A \cup (B \cup C)$, then use the distributive law
$A\cap(B\cup C) = AB \cup AC$ and apply the two-event rule again; the general case follows
by induction on $n$.
\end{proof}

\begin{example}[Two courses]
At a university, $60\%$ of students have taken Stat 230, $50\%$ have taken Math 239 and
$30\%$ have taken both. Let $A$ and $B$ be the corresponding events for a random student.
\begin{itemize}
  \item $P(\text{at least one}) = P(A\cup B) = 0.6 + 0.5 - 0.3 = 0.8$.
  \item $P(\text{neither}) = P(\bar A \bar B) = 1 - P(A\cup B) = 0.2$ (De Morgan).
  \item $P(\text{exactly one}) = P(A\cup B) - P(AB) = 0.8 - 0.3 = 0.5$.
  \item $P(\text{Stat 230 only}) = P(A\bar B) = P(A) - P(AB) = 0.3$.
\end{itemize}
\end{example}

\begin{example}[Three events]
A number is chosen uniformly at random from $\{1,2,\dots,100\}$. What is the probability it
is divisible by $2$, $3$ or $5$? Let $A_d$ be ``divisible by $d$''; then
$P(A_d) = \floor{100/d}/100$, and $A_2 \cap A_3 = A_6$ etc.
\[
  P(A_2\cup A_3\cup A_5) = \frac{50 + 33 + 20 - 16 - 10 - 6 + 3}{100} = \frac{74}{100} = 0.74.
\]
\end{example}

\begin{example}[Cards]
Draw one card from a standard $52$-card deck. Then
\[
  P(\text{heart or king}) = \frac{13}{52} + \frac{4}{52} - \frac{1}{52} = \frac{16}{52}
  = \frac{4}{13},
\]
where the $1/52$ removes the king of hearts, which was counted twice.
\end{example}

\section{Equally likely outcomes}

\begin{theorem}
If $S$ is finite and its $N$ outcomes are equally likely, then $P(a) = 1/N$ for each
outcome and
\[
  P(A) = \frac{|A|}{|S|}
\]
for every event $A$.
\end{theorem}

\begin{proof}
By finite additivity, $1 = P(S) = \sum_{a\in S} P(a) = N\cdot p$ where $p$ is the common
value, so $p = 1/N$; then $P(A) = \sum_{a\in A} 1/N = |A|/N$.
\end{proof}

So computing probabilities in equally likely models reduces to \emph{counting} the
outcomes of $A$ and $S$, which is the subject of Chapter 3. The key modelling decision is
to choose $S$ so that its outcomes really are equally likely (for example, ordered pairs
rather than sums for two dice).

\begin{example}[Three coins]
Toss three fair coins. Using the equally likely sample space
\[
  S = \{HHH,HHT,HTH,THH,HTT,THT,TTH,TTT\},
\]
\[
  P(\text{exactly two heads}) = \frac{|\{HHT,HTH,THH\}|}{8} = \frac38.
\]
If we had instead used $S' = \{0,1,2,3\}$ heads and wrongly treated it as equally likely,
we would get $1/4$, which is incorrect.
\end{example}

\subsection{Odds}

Probabilities are sometimes quoted as odds.

\begin{definition}[Odds]
The \emph{odds in favour} of an event $A$ are $P(A) : P(\bar A)$, i.e.\ the ratio
$P(A)/(1-P(A))$. The odds against $A$ are the reciprocal.
\end{definition}

\begin{example}
If the odds in favour of $A$ are $3:2$, then $P(A)/(1-P(A)) = 3/2$, so $P(A) = 3/5$. In
general, odds of $a : b$ correspond to $P(A) = a/(a+b)$.
\end{example}

\section{Common pitfalls}

\begin{itemize}
  \item $P(A\cup B) = P(A) + P(B)$ only when $A$ and $B$ are mutually exclusive. Adding
        without subtracting the overlap can produce ``probabilities'' larger than $1$.
  \item ``Mutually exclusive'' is not the same as ``independent'' (Chapter 4). In fact two
        disjoint events with positive probabilities are never independent.
  \item The complement of ``at least one'' is ``none'', not ``at most one''.
\end{itemize}

\section{Summary}

\begin{itemize}
  \item Axioms: $P(S)=1$, $P(A)\ge0$, and countable additivity over disjoint events.
  \item Derived rules: $P(\emptyset)=0$, $P(\bar A) = 1 - P(A)$, monotonicity,
        $0\le P(A) \le 1$.
  \item Inclusion--exclusion: $P(A\cup B) = P(A)+P(B)-P(AB)$; for three events add back
        the triple intersection.
  \item With $N$ equally likely outcomes, $P(A) = |A|/N$, so probability becomes counting.
  \item Odds $a:b$ in favour mean $P(A) = a/(a+b)$.
\end{itemize}
